\begin{abstract}
The YOSO design pattern mitigates adaptive corruption by having each
party prepare one message, erase its remaining state, and then send.
Is erasure essential?  Without it, the adversary can inspect every
speaker's entire state.  The apparent impossibility has a loophole:
some parties may never speak, and their secrets can remain hidden.
We call the resulting model \emph{You Atmost Speak Once} (YASO).
We study it as a deliberately restricted theoretical model for secure
multiparty computation, with a synchronous public ledger, random initial
corruptions of computation parties (helpers without application inputs
or outputs of their own),
and chosen adaptive corruptions of application parties (which supply
inputs or receive outputs), under role-specific timing rules.
We give a feasibility construction for ongoing computation with no
advance bound on the number of I/O rounds.  Independent adaptively
secure garblings exchange state through public-key ciphertexts, and
universal sampling removes the need for a secret-holding garbler.
Silent committees protect unused labels, while non-committing encryption
supports later state exposure.  The aim is to understand what remains
possible without erasure and to identify techniques for more practical
constructions.
\keywords{secure multiparty computation \and adaptive security \and
YOSO \and erasure-free protocols \and universal sampling}
\end{abstract}
